\section*{Acknowledgments}
I am deeply grateful to Seidai Yasuda for his invaluable support during the preparation of this paper.

This work is supported by JSPS KAKENHI Grant Number JP26KJ0476.



\paragraph{Use of AI.}
The author used OpenAI's ChatGPT and Codex as
discussion partners for mathematical arguments,
to assist in revising proofs and their presentation,
and to proofread the manuscript.
The author takes full responsibility for the entire content of this manuscript.


\section*{Notations and assumptions}
\label{sec:notations}
\addcontentsline{toc}{section}{Notations and assumptions}

Every category is an $(\infty,1)$-category. A space is an
$\infty$-groupoid, and $\mathcal S$ denotes the category of small spaces.

Fix nested Grothendieck universes. Small indexing categories belong to
the innermost universe in use. Presentable categories have the
corresponding small colimits, with objects in a larger universe.
Further universes are used for their categories and diagram objects,
without changing the colimit class or compactness of an existing
category.


For a category $C$, $C^\simeq$ is its maximal subgroupoid and
$\Map_C(x,y)$ its mapping space. In a stable category,
$\operatorname{map}_C(x,y)$ is its mapping spectrum.


An object $x$ is compact if $\Map(x,-)$ preserves filtered colimits.
The full subcategory of compact objects is denoted by $A^\omega$.
A presentable category $A$ is compactly generated if $A^\omega$
generates it under filtered colimits.

Write $\mathbb S$ for the sphere, $\Sp$ for the category of spectra, and
$\Sp^\omega$ for the category of finite spectra. For a small idempotent-complete
stable category $C$, $K(C)$ is connective algebraic K-theory, and
$K(R)=K(\Perf(R))$ for a ring spectrum $R$.
Ordinary Ind-completion is $\Ind$. Tensor products of presentable categories are
Lurie's \citep[Section~4.8.1]{LurieHA}.

